\section{Zeta--Schrödinger temporal branch: spectral synthesis, phase dispersion, and fuzziness}

\subsection{Why the transform must be specified}

A discrete Fourier transform/FFT assumes samples on a uniform frequency or time grid according to the chosen convention. Prime logarithms \(\log p\) and non-trivial zeta-zero ordinates are not uniformly spaced. Therefore the canonical project language is:
\begin{itemize}[leftmargin=2em]
\item uniform registry \(\Rightarrow\) DFT/FFT may be used;
\item non-uniform registry \(\Rightarrow\) direct Fourier synthesis or a declared NUFFT implementation;
\item gridding/interpolation \(\Rightarrow\) a separate approximation receipt is mandatory.
\end{itemize}
The phrase ``temporal FFT'' is therefore an implementation shortcut, not a mathematical identity between zeta and time.

\subsection{Spectral registries}

A temporal synthesis experiment begins with a content-addressed registry
\begin{equation}
\Omega=\{(\omega_j,c_j,\mathsf{source}_j)\}_{j=1}^{M}.
\end{equation}
Allowed source classes include:
\begin{enumerate}[leftmargin=2em]
\item \status{UNIFORM\_CONTROL}: a preregistered uniform frequency grid;
\item \status{PRIME\_LOG}: \(\omega_j=\Omega_0\log p_j\) for a finite prime set;
\item \status{ZETA\_ZERO\_ORDINATE}: \(\omega_j=\Omega_0\gamma_j\) for declared non-trivial zero ordinates \(\gamma_j\), used only as a diagnostic registry;
\item \status{COLLATZ\_OPERATOR\_SPECTRUM}: eigenfrequencies of a fully specified finite or mathematically controlled operator;
\item \status{EXTERNAL\_MEASUREMENT}: an experimental spectrum with provenance.
\end{enumerate}
No registry is canonical merely because it produces an attractive waveform.

\subsection{Schrödinger phase evolution}

Define a diagonal project Hamiltonian on the declared modes,
\begin{equation}
\boxed{\hat H_{ZC}|j\rangle=E_j|j\rangle,
\qquad E_j=\hbar\omega_j.}
\end{equation}
For
\begin{equation}
|\Psi(0)\rangle=\sum_j c_j|j\rangle,
\qquad\sum_j|c_j|^2=1,
\end{equation}
unitary evolution is
\begin{equation}
\boxed{
|\Psi(\tau)\rangle
=\sum_j c_j e^{-i\omega_j\tau}|j\rangle.
}
\label{eq:sch-evolution}
\end{equation}
The parameter is written \(\tau\) deliberately. It may be ordinary laboratory time in a standard Hamiltonian simulation, proper time for a physical local clock, or PNCS phase proper time in the computational substrate. Those meanings must never be silently exchanged.

A single eigenmode acquires only a global phase and has no internal relative-phase clock. At least two distinct \(\omega_j\) are required for non-trivial internal phase relations.

\subsection{Temporal synthesis}

Define the coherent temporal amplitude
\begin{equation}
\boxed{x(\tau)=\sum_{j=1}^{M}c_j e^{-i\omega_j\tau}.}
\label{eq:temporal-synthesis}
\end{equation}
For uniformly spaced frequencies this can be sampled efficiently by inverse FFT under a declared convention. For non-uniform \(\omega_j\), equation~\eqref{eq:temporal-synthesis} itself is the authority and may be evaluated directly or via NUFFT.

For finite prime-log modes,
\begin{equation}
\omega_p=\Omega_0\log p
\end{equation}
gives
\begin{equation}
\boxed{x(\tau)=\sum_{p\in P}c_p\,p^{-i\Omega_0\tau}.}
\label{eq:prime-time}
\end{equation}
Combining this with a critical-line finite prime character gives
\begin{equation}
p^{-1/2-i\Omega_0\tau}=p^{-1/2}e^{-i\Omega_0\tau\log p}.
\end{equation}
This is a finite-mode amplitude/phase identity, not use of the divergent critical-line Euler product.

\subsection{Definition of Zeta--Schrödinger fuzziness}

Let
\begin{equation}
p_j=|c_j|^2,
\qquad \sum_jp_j=1.
\end{equation}
Define the survival/coherence amplitude
\begin{equation}
\boxed{
A_Z(\tau)
=\langle\Psi(0)|\Psi(\tau)\rangle
=\sum_jp_j e^{-i\omega_j\tau}.
}
\label{eq:az}
\end{equation}
The project-defined \emph{Zeta--Schrödinger temporal fuzziness} is
\begin{equation}
\boxed{
\mathcal F_Z(\tau)=1-|A_Z(\tau)|^2.
}
\label{eq:fuzziness}
\end{equation}
Status: \status{PROJECT\_EXACT} as a definition. It satisfies
\begin{equation}
0\le \mathcal F_Z\le1,
\qquad
\mathcal F_Z(0)=0.
\end{equation}
For a one-mode state, \(\mathcal F_Z(\tau)=0\) for all \(\tau\).

\subsection{Small-time expansion}

Let
\begin{equation}
\bar\omega=\sum_jp_j\omega_j,
\qquad
(\Delta\omega)^2=\sum_jp_j(\omega_j-\bar\omega)^2.
\end{equation}
Expanding equation~\eqref{eq:az} gives
\begin{equation}
\boxed{
|A_Z(\tau)|^2
=1-(\Delta\omega)^2\tau^2+O(\tau^4),
}
\end{equation}
and hence
\begin{equation}
\boxed{
\mathcal F_Z(\tau)
=(\Delta\omega)^2\tau^2+O(\tau^4).
}
\label{eq:fuzz-small}
\end{equation}
Equivalently \(\Delta E=\hbar\Delta\omega\) controls the initial rate at which the state becomes distinguishable from its initial ray.

\subsection{Relative phase as an operational clock}

For two modes,
\begin{equation}
\Delta\varphi_{12}(\tau)
=(\omega_2-\omega_1)\tau.
\end{equation}
The measurable internal change is relative phase, not the common global phase. A project clock observable may therefore be built from a vector of pairwise phase differences or from \(\mathcal F_Z\), with revivals explicitly allowed.

\INTERP
The phrase ``fuzziness makes time flow'' is sharpened here to the following candidate statement: \emph{operational temporal change is available when a spectral superposition develops non-trivial relative phases, and \(\mathcal F_Z\) measures loss of initial phase coherence}. This is compatible with standard unitary quantum dynamics. The stronger ontological claim that physical time itself is generated by zeta/Collatz spectral fuzziness remains \status{OPEN}.

\subsection{Local proper-time binding}

A physically local form replaces the evolution parameter by proper time,
\begin{equation}
|\Psi(\tau_{\rm phys})\rangle
=\sum_jc_j e^{-i\omega_j\tau_{\rm phys}}|j\rangle.
\end{equation}
A PNCS runtime form may instead use the v0.27 content-addressed phase proper time,
\begin{equation}
|\Psi(\tau_{\rm PNCS})\rangle
=\sum_jc_j e^{-i\omega_j\tau_{\rm PNCS}}|j\rangle,
\end{equation}
with the mandatory receipt field
\begin{center}
\texttt{laboratory\_time\_modified = false}.
\end{center}
The two equations have the same numerical form but different semantics and provenance.

\subsection{Revivals are features, not failures}

Because equation~\eqref{eq:fuzziness} is generated by coherent phases, \(\mathcal F_Z\) need not increase monotonically. Finite commensurate spectra can exhibit exact revivals; generic incommensurate finite spectra can exhibit approximate recurrences. Therefore the project must not equate ``time arrow'' with monotonic fuzziness without an additional open-system/decoherence model.
